\slajdid{10-s082}
\begin{frame}[label=10-s082]
\gusttekst\setlength{\abovedisplayskip}{3pt}\setlength{\belowdisplayskip}{3pt}\setlength{\abovedisplayshortskip}{2pt}\setlength{\belowdisplayshortskip}{2pt}\begin{columns}[T,onlytextwidth]\column{.29\textwidth}\centering\input{slike/tikz/s082_ecl_diferencijalni_par.tex}\column{.69\textwidth}
Ulazni napon $V_I$ je na bazi tranzistora T1, dok je na bazi tranzistora T2 referentni, fiksan, napon $V_R$. Kao i kod drugih diferencijalnih pojačavača možemo izlaz uzeti bilo sa kolektora tranzistora T1, $V_{O1}$, bilo sa kolektora tranzistora T2, $V_{O2}$. Zbog toga ECL logička kola po pravilu imaju dva izlazna priključka.
\par\medskip Karakteristika prenosa\par\bigskip
Krenućemo od napona $V_I\ll V_R$. Po konturi $V_I$, BE T1, EB T2, $V_R$ možemo da napišemo (i to će uvek važiti)
\[V_I=V_{BE1}-V_{BE2}+V_R\]
Kako smo pretpostavili $V_I\ll V_R$
\[V_I=V_{BE1}-V_{BE2}+V_R\ll V_R\]
\[V_{BE1}-V_{BE2}\ll0\]
Uočiti da zbog strujnog izvora $I_E$ mora voditi bar jedan tranzistor.
\par\medskip (Idealan strujni izvor će prilagoditi napon na sebi tako da obezbedi putanju struju)
\end{columns}
\end{frame}
